
%(BEGIN_QUESTION)
% Copyright 2006, Tony R. Kuphaldt, released under the Creative Commons Attribution License (v 1.0)
% This means you may do almost anything with this work of mine, so long as you give me proper credit

Fra aktuatorens perspektiv: hvilken ventil er enklest å posisjonere, enkeltsetet globeventil eller dobbeltsetet globeventil (alt annet likt)? Hvorfor?

\underbar{file i00782}
%(END_QUESTION)





%(BEGIN_ANSWER)

Dobbeltsetede globeventiler er enklere å posisjonere enn enkeltsetede, fordi dobbeltsetet ventilsats reduserer kraften fluidtrykket utøver på pluggene. Med mindre prosessindusert kraft på ventilspindelen får aktuatoren lettere jobb med å flytte ventilen til ønsket posisjon.

%(END_ANSWER)





%(BEGIN_NOTES)


%INDEX% Sluttkontrollelementer, ventil: dynamisk fluidkraft på globeventilplugg
%INDEX% Sluttkontrollelementer, ventil: dobbeltsetet
%INDEX% Sluttkontrollelementer, ventil: enkeltsetet

%(END_NOTES)


